% Generated draft: requires mathematical and source review.
\documentclass[11pt]{article}
\usepackage[margin=25mm]{geometry}
\usepackage{fontspec}
\setmainfont{Latin Modern Roman}
\usepackage{amsmath,amssymb,mathtools,unicode-math}
\setmathfont{Latin Modern Math}
\usepackage{enumitem,needspace,xurl,hyperref}
\setlength{\parindent}{0pt}
\setlength{\parskip}{5pt}
\setlength{\emergencystretch}{3em}
\newcommand{\N}{\mathbb{N}}
\newcommand{\Z}{\mathbb{Z}}
\newcommand{\Q}{\mathbb{Q}}
\newcommand{\R}{\mathbb{R}}
\newcommand{\C}{\mathbb{C}}
\newcommand{\F}{\mathbb{F}}
\newcommand{\A}{\mathbb{A}}
\newcommand{\PP}{\mathbb{P}}
\newcommand{\E}{\mathbb{E}}
\newcommand{\eps}{\varepsilon}
\newcommand{\dd}{\,\mathrm{d}}
\begin{document}
\section*{UnsolvedMath ID: 30002957}
\textbf{Source problem code:} catalogue entry
\subsection{Definitions and mathematical statement}
Let \(d\geq 1\) be an integer, and let \(X\subset\R^d\) be the set of points of the Poisson point process mentioned in the source. The source does not specify its intensity measure, homogeneity, or nondegeneracy assumptions. Work on configurations for which the stated Delaunay triangulation \(T\) is defined, and regard \(T\) as an abstract simplicial complex with vertex set \(X\).

The \emph{flag complex} of the \(1\)-skeleton of \(T\) is
\[
F(T)=\bigl\{\sigma\subset X: \sigma\text{ is finite and }\{x,y\}\in T\text{ for all distinct }x,y\in\sigma\bigr\}.
\]
For each integer \(0\leq k\leq d\), write
\[
F_k(T)=\{\sigma\in F(T):|\sigma|=k+1\},\qquad
D_k(T)=\{\sigma\in F_k(T):\sigma\in T\}.
\]
Thus a flag simplex is called Delaunay precisely when it is a simplex of the given triangulation \(T\). Flag simplices are understood abstractly; their definition alone imposes no affine independence requirement.

\textbf{Source question.} For a Poisson point process in \(\R^d\), determine how many members of \(F_d(T)\), and more generally of \(F_k(T)\) for each integer \(0\leq k\leq d\), belong to \(D_k(T)\).

\textbf{Unresolved specification.} The source does not determine a unique numerical or probabilistic target. It does not specify whether ``how many'' means a random count, an expected count, a spatial intensity, or a proportion; the supplied literature assessment describes a fraction question. Counts over all of \(\R^d\) need not be finite. A finite-window formulation would require a window and a rule assigning simplices to it, while a proportion would require a sampling or limiting convention. None of these choices is imposed here. The probability law and the treatment of configurations with nonunique Delaunay triangulations also require clarification before a fully specified quantitative problem can be stated.
\subsection{Short English statement}
Start with the Delaunay triangulation of a Poisson point set in \(\R^d\). Add an abstract simplex whenever all pairs of its vertices are Delaunay edges. The question asks how many of these flag simplices are themselves Delaunay simplices, especially in dimension \(d\) and more generally for \(0\leq k\leq d\). The process law and the intended counting or fraction convention remain unspecified.
\subsection{Sources}
\begin{itemize}
\item Supplied local UnsolvedMath registry record, permanent ID \(30002957\), title ``Delaunay Simplices in Flag Complexes of Random Triangulations,'' exact source problem code \texttt{catalogue entry}; collection: Oberwolfach Reports Open Problems. The record gives proposed year 2015 and names no individual proposer.
\item Supplied source citation: \emph{Computational Geometric and Algebraic Topology} (2016), Oberwolfach Reports, pp.\ 2637--2699, DOI \texttt{10.4171/owr/2015/45}. \url{https://doi.org/10.4171/owr/2015/45}. The background describes Oberwolfach Reports as records of mathematical workshops and their open-problem discussions. The supplied literature list labels this reference ``Oberwolfach Report 45, 2015.'' The differing year labels are retained without reconciliation.
\item Supplied literature reference: Edelsbrunner--Nikitenko--Reitzner, \emph{Expected Sizes of Poisson--Delaunay Mosaics}. \url{https://doi.org/10.1017/apr.2017.4}.
\item Literature assessment checked at \texttt{2026-08-21}: the supplied record labels the problem partially solved and reports that primary stochastic-geometry work provides explicit expected simplex intensities and exact low-dimensional values for Poisson--Delaunay mosaics, but does not provide a general closed answer to the flag-complex fraction question in arbitrary dimension. This is supplied dated curation derived from the OWR material, not an independently verified assessment of the literature. Its time-sensitive status and cited claims require independent verification.
\item The supplied statement status is \texttt{corrected\_verified}, with an accompanying claim that mathematical content was preserved while notation, typography, or omitted definitions were repaired. This is registry metadata, not independent mathematical verification of this draft.
\end{itemize}
\end{document}
